\chapter{Self-improvement}
\label{chap:self-improvement}
\label{sec:self-improvement}

Throughout this chapter, $(\psi,A,B,L)$ is an $(\eps,\delta,\gamma)$-good
symmetric strategy for the $(m,q,d)$ low individual degree test.  We first
construct a sub-measurement on the original local space from a semidefinite
program.  Its dual certificate satisfies $Z^2\le H\le I$.  We then dilate any
finite family of such sub-measurements simultaneously to projective
sub-measurements on one common enlarged space.

\section{A semidefinite program for measurements}

For $g\in\polyfunc mqd$ put $A_g=\E_uA^u_{g(u)}$.  The primal and dual programs
are
\[
	\sup_{T_g\ge0,\ \sum_gT_g\le I}\ \sum_g\Tr(T_gA_g),
	\qquad
	\inf_{Z\ge A_g\ \forall g}\ \Tr(Z).
\]

\begin{lemma}[SDP certificate]\label{lem:sdp}
	\lean{MIPStarRE.LDT.SelfImprovement.matrixSdpCanonicalStrongDuality, MIPStarRE.LDT.SelfImprovement.sdp_statement_with_slackness, MIPStarRE.LDT.SelfImprovement.SdpStatementWithSlackness}
	\leanok
	\uses{def:good-strategy, def:polymeasurements}
	The two programs have the same attained optimum.  There are optimal
	solutions $T=\{T_g\}$ and $Z$ such that $T$ is a measurement, $Z\ge A_g$
	for every $g$, and $T_g(Z-A_g)=0$ for every $g$.  In particular $Z\ge0$.
\end{lemma}
\begin{proof}
	\leanok
	The block-diagonal canonical program has one block for each polynomial
	and one slack block for $I-\sum_gT_g$; its values are those of the
	displayed programs.  Both sides are strictly feasible, so
	finite-dimensional strong duality gives attained optima with complementary
	slackness.  The formalization proves strong duality by a separation
	argument for the closed image cone.  Assigning the unused part of $T$ to a
	fixed polynomial does not decrease the objective and yields a measurement.
\end{proof}

\begin{lemma}[Certificate contraction]\label{lem:sdp-certificate-contraction}
	\lean{MIPStarRE.LDT.SelfImprovement.sdp_dual_square_eq_sandwiched_average, MIPStarRE.LDT.SelfImprovement.sdp_filtered_certificate, MIPStarRE.LDT.SelfImprovement.helper_filtered_certificate_of_slackness}
	\leanok
	\uses{lem:sdp}
	Let $T,Z$ be as in Lemma~\ref{lem:sdp}, put $P^u_g=A^u_{g(u)}$, and let
	$H_g=\E_uP^u_gT_gP^u_g$ and $H=\sum_gH_g$.  Then $H$ is a sub-measurement
	and
	\[
		Z^2=\sum_gA_gT_gA_g,\qquad
		H-Z^2=\sum_g\E_u(P^u_g-A_g)T_g(P^u_g-A_g)\ge0,
	\]
	so that $Z^2\le H\le I$ and $0\le Z\le I$.
\end{lemma}
\begin{proof}
	\leanok
	Grouping $g$ by $g(u)$ gives $\sum_gP^u_gT_gP^u_g\le\sum_aA^u_a=I$.
	Complementary slackness and its adjoint give
	$Z^2=Z(\sum_gT_g)Z=\sum_gA_gT_gA_g$.  Centering $P^u_g$ at its average
	$A_g$ expresses $H-Z^2$ as an average of positive sandwiches.
\end{proof}

\section{The filtered measurement}

\begin{lemma}[Add-in-$u$]\label{lem:add-in-u}
	\lean{MIPStarRE.LDT.SelfImprovement.AddInUFullStatement, MIPStarRE.LDT.SelfImprovement.addInUFullStatement_of_isGood}
	\leanok
	\uses{lem:global-variance-of-points, lem:sdp}
	For every measurement $T$ with polynomial outcomes, the variance of the
	point measurements along the outcomes of $T$ is controlled by
	Lemma~\ref{lem:global-variance-of-points}: inserting or removing a point
	measurement $A^u_{g(u)}$ next to $T_g$ inside the expectations used below
	changes them by at most a multiple of
	$\sqrt{24m(\eps+\delta+d/q)}$.
\end{lemma}
\begin{proof}
	\leanok
	This is the add-in-$u$ lemma of \cite{Ji2020LowIndividualDegree}: a
	Cauchy--Schwarz estimate whose variance factor is the global variance of
	the point measurements, bounded by Lemma~\ref{lem:global-variance-of-points}.
\end{proof}

\begin{lemma}[Self-improvement with non-projective output]\label{lem:self-improvement-helper}
	\lean{MIPStarRE.LDT.SelfImprovement.self_improvement_helper_with_contraction, MIPStarRE.LDT.SelfImprovement.SelfImprovementHelperStatement, MIPStarRE.LDT.SelfImprovement.selfImprovementHelperError}
	\leanok
	\uses{def:good-strategy, def:polymeasurements, def:simeq, def:strong-self-consistency}
	Let $G\in\polymeas mqd$ satisfy
	$A^u_a\ot I\simeq_\nu I\ot G_{[g(u)=a]}$ on average over $u$, and put
	$\zeta_0=100m(\eps^{1/2}+\delta^{1/2}+(d/q)^{1/2})$.  Then there are
	$H\in\polysub mqd$ and a positive semidefinite $Z\le I$ with
	$Z^2\le\sum_hH_h$ such that
	\begin{enumerate}
		\item (completeness) $\bra\psi H\ot I\ket\psi\ge1-\nu-\zeta_0$;
		\item (consistency with $A$) $A^u_a\ot I\simeq_{\zeta_0}I\ot H_{[h(u)=a]}$;
		\item (strong self-consistency)
		      $\sum_h\bra\psi H_h\ot H_h\ket\psi\ge\bra\psi H\ot I\ket\psi-\zeta_0$;
		\item (boundedness)
		      $\bra\psi Z\ot I\ket\psi-\E_u\sum_a\bra\psi A^u_a\ot H_{[h(u)=a]}\ket\psi\le\zeta_0$
		      and $Z\ge\E_uA^u_{h(u)}$ for every $h$.
	\end{enumerate}
\end{lemma}
\begin{proof}
	\leanok
	\uses{lem:sdp, lem:sdp-certificate-contraction, lem:add-in-u, lem:local-variance-of-points, prop:two-notions-of-self-consistency, cor:schwartz-zippel-individual}
	Take $T$ and $Z$ from Lemma~\ref{lem:sdp} and the filtered
	sub-measurement $H$ of Lemma~\ref{lem:sdp-certificate-contraction}, which
	also gives $Z^2\le H$.  The formal proof of the four conclusions follows
	\cite{Ji2020LowIndividualDegree}.  Consistency with $A$ is the
	off-diagonal part of the add-in-$u$ estimate
	(Lemma~\ref{lem:add-in-u}).  Strong self-consistency moves the outer
	point measurements across the tensor product with point self-consistency
	and the add-in-$u$ estimate, and bounds the contribution of distinct
	polynomials agreeing at a random point by Schwartz--Zippel.  Completeness
	combines complementary slackness with the consistency of the seed $G$:
	$\bra\psi Z\ot I\ket\psi\ge\sum_g\bra\psi A_g\ot G_g\ket\psi\ge1-\nu$.
	Boundedness follows from $Z\ge A_h$ and the consistency estimate.  The
	natural error terms are absorbed into $\zeta_0$ by elementary scalar
	comparisons; if $\zeta_0\ge1$ the conclusions are immediate.
\end{proof}

\section{Simultaneous dilation and compression}
\label{sec:self-improvement-dilation}

\begin{lemma}[One-measurement Naimark dilation]\label{thm:naimark}
	\lean{MIPStarRE.LDT.MakingMeasurementsProjective.oneMeasNaimark, MIPStarRE.LDT.MakingMeasurementsProjective.OneMeasNaimarkData}
	\leanok
	\uses{def:measurements}
	Let $M=\{M_a\}_{a\in\calA}$ be a sub-measurement on $\calH$ and let
	$\calK$ have an orthonormal basis indexed by $\calA\cup\{\bot\}$.  There
	are projective measurements $P=\{P_a\}_{a\in\calA\cup\{\bot\}}$ on
	$\calH\ot\calK$ such that $J^\dagger P_aJ=M_a$ for $a\in\calA$, where
	$Jv=v\ot\ket\bot$.
\end{lemma}
\begin{proof}
	\leanok
	Complete $M$ by $M_\bot=I-\sum_aM_a$.  The map
	$V v=\sum_a\sqrt{M_a}v\ot\ket a$ is an isometry; extend the isometry
	$Jv\mapsto Vv$ to a unitary $U$ and set
	$P_a=U^\dagger(I\ot\ket a\!\bra a)U$.
\end{proof}

\begin{lemma}[Simultaneous dilation and compression]\label{lem:simultaneous-dilation-compression}
	\lean{MIPStarRE.LDT.SelfImprovement.simultaneousDilationFamily, MIPStarRE.LDT.SelfImprovement.simultaneousDilationFamily_outcome_compression, MIPStarRE.LDT.SelfImprovement.simultaneousDilationFamily_correlation, MIPStarRE.LDT.SelfImprovement.simultaneousDilationState, MIPStarRE.LDT.SelfImprovement.compressMeasurementAtNone, MIPStarRE.LDT.SelfImprovement.mixed_bipartiteConsError_of_compression, MIPStarRE.LDT.SelfImprovement.dilated_qBipartiteSSCDefect}
	\leanok
	\uses{thm:naimark, def:simeq, def:strong-self-consistency}
	Let $\{H^x\}_x$ be a finite family of sub-measurements on $\calH$ with a
	common finite outcome set.  There are a finite ancilla $\calK$, the
	isometry $Jv=v\ot\ket0$, and projective sub-measurements $P^x$ on
	$\calH\ot\calK$ with $J^\dagger P^x_gJ=H^x_g$.  For a symmetric state put
	$\ket{\psi'}=(J\ot J)\ket\psi$ and extend original operators $R$ to
	$R'=R\ot I_\calK$.  Then slice masses, consistency with original
	measurements, and correlations between one original operator and one
	slice effect are preserved, and
	\[
		\sum_g\|(P^x_g\ot I-I\ot P^x_g)\ket{\psi'}\|^2
		=2\Bigl(\bra\psi H^x\ot I\ket\psi-\sum_g\bra\psi H^x_g\ot H^x_g\ket\psi\Bigr).
	\]
	If $T$ is a measurement on $\calH\ot\calK$, then $J^\dagger T_hJ$ is a
	measurement on $\calH$ with the same consistency with the original point
	measurements.
\end{lemma}
\begin{proof}
	\leanok
	\uses{thm:naimark}
	Apply Lemma~\ref{thm:naimark} to each $H^x$ with the same ancilla.  The
	correlation identities follow from $J^\dagger P^x_gJ=H^x_g$, and the mirror
	identity by expanding the square and using projectivity and symmetry.
	Compression preserves positivity and $J^\dagger IJ=I$.
\end{proof}

\begin{theorem}[Self-improvement with dilation]\label{thm:self-improvement-in-induction-section}
	\lean{MIPStarRE.LDT.SelfImprovement.selfImprovementWithDilation, MIPStarRE.LDT.SelfImprovement.DilationSelfImprovementConclusion, MIPStarRE.LDT.SelfImprovement.selfImprovementDilationError}
	\leanok
	\uses{def:good-strategy, def:polymeasurements, def:simeq, def:approx_delta}
	Let $G\in\polymeas mqd$ satisfy $A^u_a\ot I\simeq_\nu I\ot G_{[g(u)=a]}$,
	and put $\zeta=200m(\sqrt\eps+\sqrt\delta+\sqrt{d/q})$.  There are a
	finite ancilla $\calK$, the isometry $Jv=v\ot\ket0$, and a projective
	sub-measurement $H\in\polysub mqd$ on $\calH\ot\calK$ such that, on
	$\ket{\psi'}=(J\ot J)\ket\psi$ and with $A'^u_a=A^u_a\ot I_\calK$,
	\begin{enumerate}
		\item (completeness) $\bra{\psi'}H\ot I\ket{\psi'}\ge1-\nu-\zeta$;
		\item (consistency) $A'^u_a\ot I\simeq_\zeta I\ot H_{[h(u)=a]}$;
		\item (strong self-consistency) $H_h\ot I\approx_\zeta I\ot H_h$;
		\item (boundedness) there is a positive contraction $Z$ with
		      $\bra{\psi'}Z\ot(I-H)\ket{\psi'}\le\zeta$ and
		      $Z\ge\E_uA'^u_{h(u)}$ for every $h$.
	\end{enumerate}
	The same ancilla and embedding serve every finite family of inputs, and
	extending $A,B,L$ by the identity on $\calK$ preserves projectivity and all
	test success probabilities.
\end{theorem}
\begin{proof}
	\leanok
	\uses{lem:self-improvement-helper, lem:simultaneous-dilation-compression}
	Apply Lemma~\ref{lem:self-improvement-helper} on the original space and
	dilate its output with Lemma~\ref{lem:simultaneous-dilation-compression}.
	Completeness, consistency, and boundedness transfer with error $\zeta_0$,
	and the mirror identity turns the strong self-consistency defect
	$\zeta_0$ into a squared mirror error $2\zeta_0=\zeta$.  The boundedness
	residual is bounded by dual feasibility.
\end{proof}
